%-------------------------------------------------------------------------------
%   CHAPTER 5 - CONCLUSION
%-------------------------------------------------------------------------------

\section{Дүгнэлт}

Энэхүү төгсөлтийн ажлын хүрээнд Монгол улсын жижиг, дунд бизнесийн байгууллагуудын бараа нийлүүлэлтийн үйл явцыг автоматжуулах, хялбарчлах зорилготой "Tdelivery" гар утасны аппликейшнийг амжилттай хөгжүүллээ. Судалгааны ажлын эхэнд дэвшүүлсэн зорилго, зорилтуудаа бүрэн биелүүлж, практикт ашиглах боломжтой бодит системийг бүтээсэн болно.

\subsection{Хүрсэн үр дүн}
Төслийн үр дүнд дараах үндсэн ажлуудыг гүйцэтгэж, үр дүнд хүрсэн:

\begin{enumerate}
    \item \textbf{Нэгдсэн экосистемийг бүрдүүлсэн:} Захиалагч (Дэлгүүр), Нийлүүлэгч (Бөөний төв), Хүргэгч гэсэн гурван талын оролцоог хангасан, мэдээллийн нэгдсэн урсгалтай системийг бий болгосон. Энэ нь талуудын хоорондын үл ойлголцлыг арилгаж, мэдээллийн ил тод байдлыг хангасан.
    \item \textbf{Технологийн оновчтой шийдэл:} React Native болон Expo технологийг ашиглан iOS болон Android үйлдлийн систем дээр ажиллах боломжтой кросс-платформ шийдлийг гаргасан нь хөгжүүлэлтийн зардлыг хэмнэж, зах зээлд нэвтрэх хугацааг хурдасгасан.
    \item \textbf{Serverless архитектур:} Google Firebase платформыг ашиглан найдвартай, хурдан ажиллагаатай, өргөтгөх боломжтой (scalable) backend системийг бүрдүүлсэн. Энэ нь серверийн арчилгааны зардлыг тэглэж, зөвхөн хэрэглээндээ тохирсон төлбөр төлөх боломжийг олгосон.
    \item \textbf{Бодит цагийн хяналт:} Захиалгын явцыг бодит цагт (real-time) хянах боломжийг бүрдүүлснээр хэрэглэгчдийн итгэлийг нэмэгдүүлж, "Бараа хэзээ ирэх вэ?" гэсэн асуултад хариулт өгсөн.
    \item \textbf{Цахим баримт бичиг (e-POD):} Хүргэлтийг зураг болон гарын үсгээр баталгаажуулах системийг нэвтрүүлснээр цаасан баримтын хэрэглээг бууруулж, маргаантай асуудлыг шийдвэрлэхэд хялбар болгосон.
\end{enumerate}

\subsection{Тулгарсан бэрхшээл ба шийдэл}
Системийг хөгжүүлэх явцад интернэтгүй орчинд ажиллах, өгөгдлийн синхрончлол хийх зэрэг техникийн хүндрэлүүд гарсан. Эдгээрийг шийдвэрлэхийн тулд Firebase-ийн "Offline Persistence" боломжийг ашиглаж, сүлжээнд холбогдмогц өгөгдлийг автоматаар илгээх логикийг хэрэгжүүлсэн. Мөн олон хэрэглэгч зэрэг хандах үед үүсэх ачааллыг даахын тулд Firestore-ийн индексжүүлэлтийг оновчтой хийсэн.

\subsection{Цаашдын хөгжүүлэлт}
Системийг цаашид улам сайжруулж, зах зээлд бүрэн нэвтрүүлэхийн тулд дараах ажлуудыг хийхээр төлөвлөж байна:

\begin{itemize}
    \item \textbf{Төлбөрийн нэгдсэн систем:} Банкны аппликейшн (Khan Bank, Golomt Bank) болон QPay, SocialPay зэрэг төлбөрийн хэрэгслүүдийг API-аар холбож, төлбөр тооцоог бүрэн автоматжуулах.
    \item \textbf{AI ба Data Analytics:} Захиалагчийн өмнөх худалдан авалтын түүх дээр үндэслэн бараа санал болгох (Recommendation System) болон нийлүүлэгчдэд борлуулалтын таамаглал (Sales Forecasting) гаргаж өгөх хиймэл оюун ухааны модулийг хөгжүүлэх.
    \item \textbf{Маршрут оновчлолын сайжруулалт:} Олон цэгт хүргэлт хийх үед хамгийн дөт замыг автоматаар тооцоолж, замын түгжрэлийг тооцдог алгоритмыг сайжруулах.
    \item \textbf{Вэб хувилбар:} Нийлүүлэгчдэд зориулсан удирдлагын самбар (Admin Dashboard)-ыг вэб хувилбараар хөгжүүлж, тайлан мэдээллийг том дэлгэц дээр харах боломжийг олгох.
\end{itemize}

Эцэст нь дүгнэхэд, "Tdelivery" систем нь уламжлалт худалдааны арга барилыг халж, дижитал шилжилтийг хийхэд чухал хувь нэмэр оруулах, эдийн засгийн үр өгөөжтэй, ирээдүйтэй төсөл болсон гэж үзэж байна.

\vfill